\documentclass[11pt, a4paper]{article}

% bfw-style.tex — gemeinsame Style-Definitionen für alle Handouts/Aufgabenhefte
% Voraussetzung: Kompilation mit xelatex (wegen fontspec)

\usepackage[ngerman]{babel}
\usepackage{geometry}
\geometry{a4paper, top=2.2cm, bottom=2.2cm, left=2.3cm, right=2.3cm}

\usepackage{fontspec}
% Fallback-Kette: Source Sans Pro -> Calibri -> Segoe UI -> Latin Modern Sans
% (Latin Modern ist immer mit MiKTeX vorhanden)
\IfFontExistsTF{Source Sans Pro}{%
  \setmainfont{Source Sans Pro}[Ligatures=TeX]%
}{\IfFontExistsTF{Calibri}{%
  \setmainfont{Calibri}[Ligatures=TeX]%
}{\IfFontExistsTF{Segoe UI}{%
  \setmainfont{Segoe UI}[Ligatures=TeX]%
}{%
  \setmainfont{Latin Modern Sans}[Ligatures=TeX]%
}}}
\IfFontExistsTF{Source Code Pro}{%
  \setmonofont{Source Code Pro}[Scale=0.92]%
}{\IfFontExistsTF{Consolas}{%
  \setmonofont{Consolas}[Scale=0.92]%
}{%
  \setmonofont{Latin Modern Mono}[Scale=0.92]%
}}

\usepackage{xcolor}
% Farbpalette: hell, freundlich, modern
\definecolor{bfwPrimary}{HTML}{0EA5A4}    % Petrol/Türkis - Primär
\definecolor{bfwPrimaryDark}{HTML}{0F766E} % dunkler Petrol für Headings
\definecolor{bfwAccent}{HTML}{F59E0B}     % Amber - Akzent
\definecolor{bfwSuccess}{HTML}{10B981}    % Grün - Erfolg / Lösung
\definecolor{bfwWarn}{HTML}{F97316}       % Orange - Achtung
\definecolor{bfwInk}{HTML}{0F172A}        % fast Schwarz
\definecolor{bfwInkSoft}{HTML}{475569}    % gedämpftes Grau
\definecolor{bfwBgSoft}{HTML}{F8FAFC}     % off-white
\definecolor{bfwBgTeal}{HTML}{ECFEFF}     % helles Türkis
\definecolor{bfwBgAmber}{HTML}{FEF3C7}    % helles Gelb
\definecolor{bfwBgRose}{HTML}{FFE4E6}     % helles Rosé für Eselsbrücken

\usepackage[colorlinks=true, linkcolor=bfwPrimaryDark, urlcolor=bfwPrimaryDark]{hyperref}
\usepackage{enumitem}
\setlist[itemize]{leftmargin=*, itemsep=2pt, topsep=2pt}
\setlist[enumerate]{leftmargin=*, itemsep=2pt, topsep=2pt}

\usepackage[most]{tcolorbox}
\usepackage{booktabs}
\usepackage{tikz}
\usetikzlibrary{decorations.pathmorphing, positioning, shapes.geometric, arrows.meta}

% Merkkästchen — wichtige Definition / Kernpunkt
\newtcolorbox{merksatz}[1][]{%
  enhanced, breakable,
  colback=bfwBgTeal,
  colframe=bfwPrimary,
  boxrule=0.6pt,
  arc=4pt,
  left=10pt, right=10pt, top=8pt, bottom=8pt,
  fonttitle=\bfseries\color{bfwPrimaryDark},
  title={\faLightbulb\ Merksatz},
  #1
}

% Eselsbrücke — Mnemoniks
\newtcolorbox{eselsbruecke}[1][]{%
  enhanced, breakable,
  colback=bfwBgRose,
  colframe=bfwAccent,
  boxrule=0.6pt,
  arc=4pt,
  left=10pt, right=10pt, top=8pt, bottom=8pt,
  fonttitle=\bfseries\color{bfwWarn},
  title={Eselsbrücke},
  #1
}

% Beispiel-Box
\newtcolorbox{beispiel}[1][]{%
  enhanced, breakable,
  colback=bfwBgSoft,
  colframe=bfwInkSoft,
  boxrule=0.4pt,
  arc=3pt,
  left=10pt, right=10pt, top=8pt, bottom=8pt,
  fonttitle=\bfseries\color{bfwInk},
  title={Beispiel},
  #1
}

% Lösungs-Box (für Aufgabenhefte)
\newtcolorbox{loesung}[1][]{%
  enhanced, breakable,
  colback=bfwBgAmber!40,
  colframe=bfwSuccess,
  boxrule=0.6pt,
  arc=3pt,
  left=10pt, right=10pt, top=8pt, bottom=8pt,
  fonttitle=\bfseries\color{bfwSuccess!70!black},
  title={Lösung},
  #1
}

% Glossar-Eintrag
\newcommand{\glossareintrag}[2]{%
  \par\noindent\textbf{\color{bfwPrimaryDark}#1} \enspace #2\par\smallskip
}

% Section-Stil — etwas farbiger als Standard
\usepackage{titlesec}
\titleformat{\section}
  {\Large\bfseries\color{bfwPrimaryDark}}
  {\thesection}{0.6em}{}
\titleformat{\subsection}
  {\large\bfseries\color{bfwPrimary}}
  {\thesubsection}{0.5em}{}
\titleformat{\subsubsection}
  {\normalsize\bfseries\color{bfwInk}}
  {\thesubsubsection}{0.4em}{}

% TikZ-Stil "skizzenhaft" — etwas weicher als Rough.js, aber stabil
\tikzset{
  roughbox/.style = {
    draw=bfwPrimaryDark, thick, fill=bfwBgTeal,
    rounded corners=4pt, inner sep=6pt
  },
  roughaccent/.style = {
    draw=bfwAccent, thick, fill=bfwBgAmber,
    rounded corners=4pt, inner sep=6pt
  },
  roughline/.style = {
    draw=bfwInkSoft, thick, -{Stealth[length=2.5mm]}
  }
}

% Bullet-Symbole (bewusst kein FontAwesome — vermeidet Font-Installations-Probleme)
\newcommand{\faLightbulb}{$\bullet$}
\newcommand{\faBookmark}{$\bullet$}

% Header/Footer
\usepackage{fancyhdr}
\pagestyle{fancy}
\fancyhf{}
\renewcommand{\headrulewidth}{0pt}
\fancyfoot[L]{\small\color{bfwInkSoft}\thepage}
\fancyfoot[R]{\small\color{bfwInkSoft} BFW M-064 Beschaffung im Büromanagement}


\title{\vspace*{-1.5cm}\Huge\bfseries\color{bfwPrimaryDark}Aufgabenheft Tag~6\\[0.4em]
  \large\color{bfwInkSoft}\textnormal{Kaufvertragsstörungen --- Übungen im IHK-Klausurformat}}
\author{\textsf{Beschaffung im Büromanagement}\\
\textsf{\small\copyright\ Can Siebert\,\textperiodcentered\,2026}}
\date{Selbstlernmaterial \textperiodcentered{} 19.05.2026}

\begin{document}
\maketitle
\thispagestyle{empty}

\begin{merksatz}[title={\bfseries So nutzen Sie dieses Aufgabenheft}]
Bearbeiten Sie alle Aufgaben zunächst \emph{ohne} in die Lösungen zu schauen.
Beachten Sie die \emph{Operatoren} und schreiben Sie Antworten in vollständigen Sätzen.
Lösungen ab Seite~\pageref{loesungen}.
\end{merksatz}

\tableofcontents
\vspace{1em}
\hrule

\section{Aufgaben}

\subsection*{Aufgabe~1 --- Überblick Leistungsstörungen (8 Punkte)}

\noindent\textbf{\color{bfwAccent}YouTube-Vorbereitung:}\enspace \href{https://www.youtube.com/watch?v=JzmjVD8Bviw}{Die Kaufvertragsstörungen einfach erklärt} (\url{https://www.youtube.com/watch?v=JzmjVD8Bviw})

\medskip
\begin{enumerate}[label=(\alph*)]
  \item \textbf{Nennen} Sie die drei Hauptarten von Leistungsstörungen beim Kaufvertrag. \hfill (3~P.)
  \item \textbf{Erläutern} Sie kurz, was eine ,,Leistungsstörung'' rechtlich bedeutet. \hfill (2~P.)
  \item \textbf{Ordnen} Sie folgende Fälle der richtigen Störungsart zu (Lieferungsverzug, Annahmeverzug, Schlechtleistung): Lieferant liefert das vereinbarte Datum nicht ein \textbar{} Käufer verweigert grundlos die Annahme \textbar{} gelieferte Aktenordner sind beschädigt. \hfill (3~P.)
\end{enumerate}

\subsection*{Aufgabe~2 --- Lieferungsverzug und Verzugszinsen (12 Punkte)}

\noindent\textbf{\color{bfwAccent}YouTube-Vorbereitung:}\enspace \href{https://www.youtube.com/watch?v=i5b0MaVaKHY}{Lieferungsverzug --- Voraussetzungen (Teil~1)} (\url{https://www.youtube.com/watch?v=i5b0MaVaKHY})

\medskip
Die \emph{Bürofix AG} hat am 02.\,Mai 2026 eine Lieferung von 50~Druckern bestellt.
Vereinbart war Lieferung bis spätestens 15.\,Mai 2026. Bis zum 25.\,Mai ist nichts
eingetroffen.

\begin{enumerate}[label=(\alph*)]
  \item \textbf{Nennen} Sie die drei Voraussetzungen des Schuldnerverzugs nach §\,286 BGB. \hfill (3~P.)
  \item \textbf{Erläutern} Sie, wann eine Mahnung \emph{entbehrlich} ist (§\,286 Abs.\,2 BGB). \hfill (3~P.)
  \item \textbf{Berechnen} Sie die Verzugszinsen für eine offene Rechnung von 8.000~€ netto bei 30~Tagen Verzug (B2B-Zinssatz: Basiszins 3{,}62\,\% $+$ 9\,\%-Punkte $=$ 12{,}62\,\%, kaufmännisches Jahr $=$ 360~Tage). \hfill (3~P.)
  \item \textbf{Beurteilen} Sie: Ist der Lieferant der Bürofix AG im Verzug? Begründen Sie. \hfill (3~P.)
\end{enumerate}

\subsection*{Aufgabe~3 --- Annahmeverzug (8 Punkte)}

\begin{enumerate}[label=(\alph*)]
  \item \textbf{Erläutern} Sie den Begriff \emph{Annahmeverzug} nach §\,293 BGB. \hfill (3~P.)
  \item \textbf{Nennen} Sie zwei Rechtsfolgen des Annahmeverzugs für den Gläubiger. \hfill (2~P.)
  \item \textbf{Beurteilen} Sie folgenden Fall: Die \emph{Schreibwaren-Vertrieb GmbH} liefert vereinbarungsgemäß 200~Aktenordner. Die Käuferin verweigert ohne Begründung die Annahme. Welche Rechte hat der Lieferant? \hfill (3~P.)
\end{enumerate}

\subsection*{Aufgabe~4 --- Schlechtleistung und Käuferrechte (16 Punkte)}

\noindent\textbf{\color{bfwAccent}YouTube-Vorbereitung:}\enspace \href{https://www.youtube.com/watch?v=5GiX_pRGjdg}{§\,437 BGB --- Rechte des Käufers bei Mängeln} (\url{https://www.youtube.com/watch?v=5GiX_pRGjdg})

\medskip
Sie sind kaufmännische Sachbearbeiterin der \emph{Müller GmbH}. Sie haben 10~Bürostühle
bestellt. Bei der Lieferung am 18.\,Mai 2026 stellen Sie fest, dass die Hydraulik aller
Stühle defekt ist.

\begin{enumerate}[label=(\alph*)]
  \item \textbf{Nennen} Sie die vier Käuferrechte nach §\,437 BGB. \hfill (4~P.)
  \item \textbf{Erläutern} Sie den \emph{Vorrang der Nacherfüllung} und die zwei Nacherfüllungsformen (§\,439 BGB). \hfill (4~P.)
  \item \textbf{Beschreiben} Sie, wann der Käufer vom Vertrag \emph{zurücktreten} darf (§\,323 BGB). \hfill (3~P.)
  \item \textbf{Verfassen} Sie eine kurze Mängelanzeige an den Lieferanten \emph{Office-Möbel KG} mit konkretem Nacherfüllungsverlangen und einer Frist von 14~Tagen. \hfill (5~P.)
\end{enumerate}

\subsection*{Aufgabe~5 --- Minderung berechnen (8 Punkte)}

Ein Drucker wurde für 600~€ netto gekauft. Er ist mangelhaft (Druckqualität deutlich
schlechter als zugesichert), eine Nacherfüllung ist gescheitert. Marktwert mangelfrei:
600~€, mangelhaft: 420~€.

\begin{enumerate}[label=(\alph*)]
  \item \textbf{Nennen} Sie die Formel zur Berechnung der Minderung nach §\,441 Abs.\,3 BGB. \hfill (2~P.)
  \item \textbf{Berechnen} Sie den geminderten Kaufpreis. \hfill (3~P.)
  \item \textbf{Beurteilen} Sie: Welche Vorteile bietet die Minderung gegenüber dem Rücktritt für den Käufer? \hfill (3~P.)
\end{enumerate}

\subsection*{Aufgabe~6 --- Besonderheiten der Kaufvertragsarten (8 Punkte)}

\begin{enumerate}[label=(\alph*)]
  \item \textbf{Unterscheiden} Sie Handelskauf und Verbrauchsgüterkauf in jeweils einem Satz. \hfill (2~P.)
  \item \textbf{Erläutern} Sie die \emph{Beweislastumkehr} beim Verbrauchsgüterkauf (§\,477 BGB). \hfill (3~P.)
  \item \textbf{Unterscheiden} Sie Stückkauf und Gattungskauf --- mit je einem Beispiel aus dem Bürokontext. \hfill (3~P.)
\end{enumerate}

\vspace{1em}
\noindent\textbf{Punkte gesamt: 60}

\newpage

\section{Lösungen}\label{loesungen}

\subsection*{Lösung Aufgabe~1}

\begin{loesung}
\textbf{(a)} Drei Hauptarten: \emph{Schuldnerverzug} (Lieferungsverzug), \emph{Gläubigerverzug} (Annahmeverzug), \emph{Schlechtleistung} (mangelhafte Lieferung).

\textbf{(b)} Eine Leistungsstörung liegt vor, wenn die geschuldete Leistung nicht, nicht rechtzeitig oder nicht ordnungsgemäß erbracht wird.

\textbf{(c)} Zuordnung:
\begin{itemize}[itemsep=0pt]
  \item Lieferant liefert nicht $\to$ \textbf{Lieferungsverzug} (Schuldnerverzug).
  \item Käufer verweigert Annahme $\to$ \textbf{Annahmeverzug} (Gläubigerverzug).
  \item Aktenordner sind beschädigt $\to$ \textbf{Schlechtleistung}.
\end{itemize}
\end{loesung}

\subsection*{Lösung Aufgabe~2}

\begin{loesung}
\textbf{(a)} Voraussetzungen Schuldnerverzug nach §\,286 BGB:
\begin{enumerate}[itemsep=0pt]
  \item Fälligkeit der Leistung,
  \item Mahnung durch den Gläubiger,
  \item Vertretenmüssen des Schuldners (Vorsatz/Fahrlässigkeit, § 276 BGB).
\end{enumerate}

\textbf{(b)} Mahnung entbehrlich, wenn:
\begin{itemize}[itemsep=0pt]
  \item ein \emph{kalendermäßig bestimmter} Termin vereinbart ist,
  \item ein Termin sich \emph{kalendermäßig berechnen} lässt (z.\,B.\ ,,14 Tage nach Bestellung''),
  \item der Schuldner die Leistung \emph{ernsthaft und endgültig} verweigert,
  \item bei Geldschulden 30 Tage nach Fälligkeit und Rechnungszugang (§ 286 III BGB).
\end{itemize}

\textbf{(c)} Verzugszinsen $= 8\,000 \times 0{,}1262 \times \frac{30}{360} = \mathbf{84{,}13~€}$.

\textbf{(d)} Ja, der Lieferant ist im Verzug. Es liegt ein \emph{kalendermäßig bestimmter}
Liefertermin (15.\,Mai) vor --- die Mahnung ist nach §\,286 Abs.\,2 Nr.\,1 BGB entbehrlich.
Mit Ablauf des 15.\,Mai befindet sich der Lieferant automatisch in Verzug, sofern er
die Verspätung zu vertreten hat (was vermutet wird).
\end{loesung}

\subsection*{Lösung Aufgabe~3}

\begin{loesung}
\textbf{(a)} Annahmeverzug nach §\,293 BGB liegt vor, wenn der Gläubiger die ihm
ordnungsgemäß angebotene Leistung nicht annimmt oder eine erforderliche Mitwirkung
unterlässt.

\textbf{(b)} Rechtsfolgen für den Schuldner: \emph{Haftungsmilderung} (§ 300 BGB ---
nur noch Vorsatz und grobe Fahrlässigkeit), Anspruch auf \emph{Mehraufwendungen}
(§ 304 BGB --- Lager-, Transportkosten), Gefahrübergang (§ 446 BGB iVm § 300 BGB).

\textbf{(c)} Rechte des Lieferanten: Anspruch auf Ersatz der Mehraufwendungen
(Lagerung, Lieferversuch); ggf.\ Hinterlegung der Ware (§ 372 BGB) oder Selbsthilfe-
verkauf nach § 373 HGB. Die Käuferin bleibt zur Zahlung verpflichtet, da die Leistung
ordnungsgemäß angeboten wurde.
\end{loesung}

\subsection*{Lösung Aufgabe~4}

\begin{loesung}
\textbf{(a)} Vier Käuferrechte nach § 437 BGB:
\begin{enumerate}[itemsep=0pt]
  \item Nacherfüllung (§ 439 BGB) --- Vorrang!
  \item Rücktritt (§§ 437 Nr.\,2, 323 BGB),
  \item Minderung (§ 441 BGB),
  \item Schadensersatz (§§ 437 Nr.\,3, 280, 281 BGB).
\end{enumerate}

\textbf{(b)} \emph{Vorrang der Nacherfüllung}: Der Käufer muss dem Verkäufer zunächst
Gelegenheit zur Mängelbeseitigung geben. Erst danach sind die übrigen Rechte möglich.
Zwei Formen: \emph{Nachbesserung} (Reparatur) und \emph{Ersatzlieferung} (mangelfreie Sache).
Wahlrecht hat der Käufer.

\textbf{(c)} Rücktritt nach § 323 BGB ist möglich, wenn der Käufer dem Verkäufer
\emph{eine angemessene Frist zur Nacherfüllung} gesetzt hat und diese erfolglos
abgelaufen ist. Bei unerheblichen Mängeln ist der Rücktritt ausgeschlossen
(§ 323 V 2 BGB).

\textbf{(d)} \emph{Mängelanzeige (Beispiel)}:

\medskip
\textit{An die Office-Möbel KG\\
Sehr geehrte Damen und Herren,\\
mit Lieferschein-Nr.\ 2026-0518 vom 18.\,Mai 2026 erhielten wir 10~Bürostühle. Bei der
Wareneingangskontrolle stellten wir fest, dass die Hydraulik aller Stühle defekt ist
--- die Sitzhöhe lässt sich nicht arretieren.\\
Wir verlangen Nacherfüllung in Form der \emph{Ersatzlieferung}. Bitte liefern Sie bis
spätestens 02.\,Juni 2026 zehn mangelfreie Bürostühle gleicher Art. Die mangelhafte
Ware halten wir für Sie zur Abholung bereit.\\
Sollte die Ersatzlieferung erfolglos bleiben, behalten wir uns vor, vom Vertrag
zurückzutreten oder den Kaufpreis zu mindern, sowie Schadensersatz geltend zu machen.\\
Mit freundlichen Grüßen \textbar{} Müller GmbH \textbar{} \today}
\end{loesung}

\subsection*{Lösung Aufgabe~5}

\begin{loesung}
\textbf{(a)} Formel: \emph{Geminderter Preis} $= \text{Kaufpreis} \times \dfrac{\text{Wert mangelhaft}}{\text{Wert mangelfrei}}$.

\textbf{(b)} Geminderter Preis $= 600 \times \frac{420}{600} = \mathbf{420~€}$.
Der Käufer kann 180~€ vom bereits gezahlten Kaufpreis zurückfordern.

\textbf{(c)} Vorteile der Minderung gegenüber dem Rücktritt:
\begin{itemize}[itemsep=0pt]
  \item Der Käufer behält die Ware (oft schnell brauchbar trotz Mangel).
  \item Keine Rückabwicklung erforderlich --- spart Aufwand.
  \item Auch bei \emph{unerheblichen} Mängeln möglich (Rücktritt nicht).
\end{itemize}
\end{loesung}

\subsection*{Lösung Aufgabe~6}

\begin{loesung}
\textbf{(a)} \emph{Handelskauf}: Kaufvertrag zwischen zwei Kaufleuten (B2B) --- es gelten
zusätzlich §§ 373\,ff.\ HGB. \emph{Verbrauchsgüterkauf}: Kauf einer beweglichen Sache
durch einen Verbraucher von einem Unternehmer (§ 474 BGB) --- mit Schutzregeln zugunsten
des Käufers.

\textbf{(b)} Beweislastumkehr (§ 477 BGB): In den ersten 12~Monaten nach Übergabe wird
\emph{vermutet}, dass ein auftretender Sachmangel bereits bei Übergabe vorhanden war.
Der Verkäufer muss das Gegenteil beweisen. Diese Regel schützt den Verbraucher, der
sonst die schwierige Aufgabe hätte, die ursprüngliche Mangelhaftigkeit nachzuweisen.

\textbf{(c)} \emph{Stückkauf}: konkrete, individuell bestimmte Sache (z.\,B.\ ein
gebrauchter Drucker mit Seriennummer 12345). \emph{Gattungskauf}: nach Gattung bestimmte
Sache (z.\,B.\ ,,50 Tonerkartuschen HP~26A''). Beim Gattungskauf kann der Lieferant
aus seinem Lagerbestand eine andere Sache derselben Gattung liefern; beim Stückkauf
wird die Leistung unmöglich, wenn die einzige Sache untergeht.
\end{loesung}

\vfill
\hrule
\vspace{0.5em}
\noindent\small\color{bfwInkSoft}\copyright\ Can Siebert\,\textperiodcentered\,2026\,\textperiodcentered{} Beschaffung im Büromanagement\,\textperiodcentered{} Tag~6 Aufgabenheft

\end{document}
